% ---- pg-020 （书页 8） ----

\[
\bar{\bm{x}} = \begin{bmatrix} \bar{x}_1 \\ \bar{x}_2 \end{bmatrix} = \begin{bmatrix} 50 \\ 4 \end{bmatrix}
\]

样本方差和协方差是

\[
\begin{aligned}
s_{11} &= \frac{1}{4}\sum_{j=1}^{4}(x_{j1}-\bar{x}_1)^2 \\
&= \frac{1}{4}\times[(42-50)^2+(52-50)^2+(48-50)^2+(58-50)^2]=34
\end{aligned}
\]

\[
\begin{aligned}
s_{22} &= \frac{1}{4}\sum_{j=1}^{4}(x_{j2}-\bar{x}_2)^2 \\
&= \frac{1}{4}\times[(4-4)^2+(5-4)^2+(4-4)^2+(3-4)^2]=0.5
\end{aligned}
\]

\[
\begin{aligned}
s_{12} &= \frac{1}{4}\sum_{j=1}^{4}(x_{j1}-\bar{x}_1)(x_{j2}-\bar{x}_2) \\
&= \frac{1}{4}\times[(42-50)(4-4)+(52-50)(5-4) \\
&\quad +(48-50)(4-4)+(58-50)(3-4)]=-1.5
\end{aligned}
\]

和

\[
s_{21}=s_{12}
\]

\[
\bm{S}_n = \begin{bmatrix} 34 & -1.5 \\ -1.5 & 0.5 \end{bmatrix}
\]

样本相关系数是

\[
r_{12} = \frac{s_{12}}{\sqrt{s_{11}}\sqrt{s_{22}}} = \frac{-1.5}{\sqrt{34}\times\sqrt{0.5}} = -0.36
\]

\[
r_{21}=r_{12}
\]

故

\[
\bm{R} = \begin{bmatrix} 1 & -0.36 \\ -0.36 & 1 \end{bmatrix}
\]

\mvsub{图解法}

作图是进行数据分析的重要的辅助手段,但常常被人们忽略.尽管不可能同时对几个变量的所有测量值作图并研究它们的形状,但单个变量的图形和一对变量的图形仍可提供很多信息.高级计算机程序和显示设备使我们可以相对容易地从一维、二维或三维空间轻松而直观地检查数据.另一方面,利用纸和笔绘图,也可从数据中领悟许多有价值的东西.表示数据的简单、优美且有效的方法可以从文献[29]得到.为一对对的变量作图并直观地考察连结的模式,这是一种很好的统计学做法.考虑以下两个变量的 7 对观测值:

变量 1($x_1$): \quad 3 \quad 4 \quad 2 \quad 6 \quad 8 \quad 2 \quad 5

变量 2($x_2$): \quad 5 \quad 5.5 \quad 4 \quad 7 \quad 10 \quad 5 \quad 7.5

这些数据在二维平面上用 7 个点表示(每个坐标轴代表一个变量),如\hyperref[fig:1.1]{图 1.1} 所示.每个点的坐标由一对测量值确定: (3,5),(4,5.5),$\cdots$,(5,7.5),所得到的二维图形称为散布图表或散布图.

\hyperref[fig:1.1]{图 1.1} 还分别显示了由变量 1 和变量 2 的观测值单独作的散点图.这种图称为(边缘)点%
